% ============================================================================
% CAPÍTULO 17 — Docker\index{Docker} e contêineres
% Status: texto completo (rodada editorial 2)
% ============================================================================
\chapter{Docker e Contêineres}
\label{cap:docker}

\begin{caixaconceito}
\textbf{Docker} empacota sua aplicação com tudo o que ela precisa (runtime,
dependências, config) em um \emph{contêiner} — que roda igual em qualquer
máquina. É o padrão da indústria \cite{docker-docs} e o pré-requisito para
CI/CD e cloud (capítulos~\ref{cap:cicd} e~\ref{cap:cloud}).
\end{caixaconceito}

Nos capítulos anteriores, você montou um projeto que depende de PostgreSQL,
Redis e de um runtime específico de Node.js. Agora pense no que acontece quando
esse projeto precisa rodar na máquina de um colega, no servidor de CI e em
produção: três ambientes, três versões de dependências, três formas de
quebrar. Este capítulo apresenta a ferramenta que elimina esse problema pela
raiz.

Ao final deste capítulo, você será capaz de:
\begin{objetivos}
  \item Explicar contêineres vs máquinas virtuais e o papel das imagens;
  \item Escrever \texttt{Dockerfile} otimizados (multi-stage\index{multi-stage}, camadas, usuário não-root\index{não-root});
  \item Orquestrar serviços com \texttt{docker-compose} (app + banco + redis);
  \item Usar volumes\index{volumes}, redes, health checks e variáveis de ambiente;
  \item Publicar imagens em um registry;
  \item Entregar o projeto do capítulo: o ambiente de desenvolvimento completo.
\end{objetivos}

\section{O problema que os contêineres resolvem}

``Funciona na minha máquina'' é o maior custo invisível do desenvolvimento.
Contêineres empacotam \emph{aplicativo + ambiente}: a mesma imagem\index{imagem} roda no seu
notebook, no CI e na nuvem. Ao contrário de uma VM (sistema operacional
completo), um contêiner\index{contêiner} compartilha o kernel do host — por isso é leve e
inicia em segundos.

\begin{caixaconceito}
\textbf{Imagem} = modelo imutável (receita). \textbf{Contêiner} = instância em
execução (prato servido). \textbf{Volume} = armazenamento persistente que
sobrevive ao contêiner. \textbf{Rede} = como os contêineres conversam entre si.
\end{caixaconceito}

\subsection{Como um contêiner isola de verdade}

Um contêiner não é uma máquina virtual com um sistema operacional próprio. Ele
é um \emph{processo comum do host} que roda dentro de um ambiente isolado
construído com dois mecanismos do kernel Linux:

\begin{itemize}[leftmargin=*, itemsep=0.4em]
  \item \textbf{Namespaces} — cada contêiner enxerga sua própria visão do
  sistema: PID, rede, sistema de arquivos, usuários. O processo dentro do
  contêiner ``acha'' que é o processo 1 daquela máquina;
  \item \textbf{Cgroups} — o kernel limita quanto de CPU, memória e I/O cada
  contêiner pode usar, impedindo que um serviço consuma o host inteiro.
\end{itemize}

Isso explica as propriedades práticas que você já deve ter ouvido: contêineres
iniciam em segundos (não há boot de SO), são leves (compartilham o kernel) e
são ``descartáveis'' (criar e destruir é barato). No Windows/macOS, o Docker
roda uma VM mínima com kernel Linux para fornecer esses mesmos mecanismos —
você não precisa se preocupar com isso, mas é útil entender por que ``funciona
igual em qualquer máquina''.

\subsection{Imagens e camadas}

Uma imagem é construída por \emph{camadas} (\emph{layers}): cada instrução do
\texttt{Dockerfile} cria uma camada imutável sobre a anterior. Duas
consequências práticas:

\begin{itemize}[leftmargin=*, itemsep=0.4em]
  \item \textbf{Cache}: se uma camada não mudou, o Docker a reaproveita do
  cache de build. Por isso a ordem importa: copie \texttt{package.json} antes
  do código — as dependências só são reinstaladas quando elas mudam;
  \item \textbf{Compartilhamento}: camadas iguais entre imagens são baixadas
  uma única vez. O Docker Hub armazena e distribui por camada.
\end{itemize}

\section{Escrevendo um Dockerfile de qualidade}

\begin{lstlisting}[style=livro, language=Bash,
  caption={Dockerfile multi-stage para uma app Next.js.}, label={list:docker-dockerfile}]
# ---------- Estágio 1: build ----------
FROM node:24-alpine AS build
WORKDIR /app
COPY package*.json ./
RUN npm ci
COPY . .
RUN npm run build

# ---------- Estágio 2: produção (mínimo) ----------
FROM node:24-alpine AS run
WORKDIR /app
ENV NODE_ENV=production
RUN addgroup -S app && adduser -S app -G app   # usuário não-root
COPY --from=build /app/.next ./.next
COPY --from=build /app/public ./public
COPY --from=build /app/package.json ./package.json
COPY --from=build /app/node_modules ./node_modules
USER app
EXPOSE 3000
CMD ["npm", "start"]
\end{lstlisting}

\begin{caixadica}
Boas práticas: (1) \textbf{multi-stage} — a imagem final só tem o que roda;
(2) \textbf{ordem das camadas} — \texttt{package.json} antes do código para
reaproveitar o cache de \texttt{npm ci}; (3) \textbf{usuário não-root} —
princípio do menor privilégio; (4) imagens \emph{pinadas} (evite \texttt{latest}).
\end{caixadica}

\begin{caixaconceito}
\textbf{Multi-stage} é o recurso que resolve o dilema ``preciso das
ferramentas de build, mas não quero elas na imagem final''. O primeiro estágio
compila (TypeScript, Next.js, Prisma client); o segundo copia apenas os
artefatos prontos. O resultado: uma imagem pequena, com menos superfície de
ataque e menos para atualizar.
\end{caixaconceito}

\begin{caixadica}
O Next.js e o Prisma precisam de atenção extra no Docker: gere o Prisma Client
no estágio de build e copie a pasta \texttt{node\_modules/.prisma}; para o
Next.js \emph{standalone}, use \texttt{output: "standalone"} no
\texttt{next.config.ts} e copie a pasta \texttt{.next/standalone} — a imagem
cai para uma fração do tamanho.
\end{caixadica}

\section{Docker Compose: o ambiente de desenvolvimento}

O Dockerfile\index{Dockerfile} descreve \emph{uma} imagem. Um produto full stack tem vários
serviços — banco, cache, fila, a própria aplicação. O Compose declara todos
eles em um arquivo YAML e os orquestra com um único comando:

\begin{lstlisting}[style=livro, language=Bash,
  caption={docker-compose.yml com app, banco e redis.}, label={list:docker-compose}]
services:
  app:
    build: .
    ports:
      - "3000:3000"
    environment:
      DATABASE_URL: postgres://app:senha@db:5432/app
      REDIS_URL: redis://redis:6379
    depends_on:
      db:
        condition: service_healthy
      redis:
        condition: service_healthy

  db:
    image: postgres:18
    environment:
      POSTGRES_USER: app
      POSTGRES_PASSWORD: senha
      POSTGRES_DB: app
    volumes:
      - dados_postgres:/var/lib/postgresql/data
    healthcheck:
      test: ["CMD-SHELL", "pg_isready -U app"]
      interval: 5s
      timeout: 5s
      retries: 10

  redis:
    image: redis:7-alpine
    healthcheck:
      test: ["CMD", "redis-cli", "ping"]
      interval: 5s
      timeout: 3s
      retries: 5

volumes:
  dados_postgres:
\end{lstlisting}

\begin{lstlisting}[style=livro, language=Bash,
  caption={Comandos essenciais do Compose.}, label={list:docker-comandos}]
docker compose up -d          # sobe tudo em segundo plano
docker compose up --build     # reconstrói a imagem do app
docker compose logs -f app    # acompanha os logs
docker compose down           # derruba (mantém volumes)
docker compose down -v        # derruba e apaga volumes (cuidado!)
\end{lstlisting}

\begin{caixadica}
Note as variáveis de ambiente: \texttt{DATABASE\_URL} aponta para
\texttt{db}, não \texttt{localhost}. Dentro da rede do Compose, os serviços
se encontram pelo \emph{nome do serviço} — o Docker resolve
\texttt{db:5432} para o contêiner do PostgreSQL. Essa é uma das diferenças
mais comuns entre ``funciona no meu computador'' e ``funciona no Docker''.
\end{caixadica}

\section{Health checks, redes e depuração}

\begin{itemize}[leftmargin=*, itemsep=0.4em]
  \item \textbf{Health check}: o Compose só inicia dependências saudáveis (\texttt{depends\_on: condition: service\_healthy}) — elimina o ``banco ainda não subiu'';
  \item \textbf{Redes}: serviços conversam pelo nome do serviço (\texttt{db}, \texttt{redis});
  \item \textbf{Depuração}: \texttt{docker exec -it app sh}, \texttt{docker ps}, \texttt{docker logs};
  \item \textbf{Produção}: imagens pequenas, sem secrets na imagem, registry privado + scan de vulnerabilidades.
\end{itemize}

\begin{caixadica}
\texttt{depends\_on} sem health check apenas ordena a criação — o PostgreSQL
pode ``existir'' mas ainda não aceitar conexões. Com
\texttt{condition: service\_healthy}, o Compose espera o banco responder
\texttt{pg\_isready} antes de subir o app. É a diferença entre ``às vezes
funciona'' e ``sempre funciona''.
\end{caixadica}

\section{Publicando imagens em um registry}

Para rodar a imagem em outro lugar — CI, servidor, nuvem — ela precisa ir para
um \emph{registry}. O fluxo:

\begin{enumerate}[leftmargin=*, itemsep=0.4em]
  \item \texttt{docker build -t seu-usuario/skillhub:1.0.0 .} — constrói e etiqueta;
  \item \texttt{docker push seu-usuario/skillhub:1.0.0} — publica no registry;
  \item \texttt{docker pull seu-usuario/skillhub:1.0.0} — baixa em qualquer máquina;
  \item \texttt{docker run -p 3000:3000 seu-usuario/skillhub:1.0.0} — executa.
\end{enumerate}

\begin{caixaarmadilha}
A imagem é pública para quem tem acesso ao registry. Nunca grave
\texttt{.env}, tokens ou senhas na imagem — injete por variável de ambiente ou
sistema de segredos no momento da execução. Um dump de uma imagem antiga não
pode vazar nada além do código.
\end{caixaarmadilha}

% ============================================================================
\section{Projeto do capítulo: ambiente de desenvolvimento completo}
\label{sec:projeto17}

\begin{caixaprojeto}
\textbf{Objetivo:} dockerizar o SkillHub (capítulo~\ref{cap:fullstack}): app
Next.js + PostgreSQL + Redis em um único \texttt{docker compose up}.

\textbf{Critérios de aceite:}
\begin{itemize}[leftmargin=*, itemsep=0.2em]
  \item Dockerfile multi-stage com usuário não-root e build reproduzível;
  \item \texttt{docker-compose.yml} com health checks e volumes persistentes;
  \item Migrações e seed rodando automaticamente na primeira subida;
  \item Imagem de produção $<$ 300 MB e sem secrets;
  \item \texttt{docker compose up -d} sobe o app 100\% funcional em uma máquina limpa;
  \item \texttt{.dockerignore} correto (sem \texttt{node\_modules}, \texttt{.next}, \texttt{.env}).
\end{itemize}
\end{caixaprojeto}

\subsection{Passo a passo}

\begin{passos}
  \item \textbf{Dockerfile}: escreva o multi-stage (build + produção) com
  usuário não-root e \texttt{output: "standalone"};
  \item \textbf{.dockerignore}: exclua \texttt{node\_modules}, \texttt{.next},
  \texttt{.env*}, \texttt{.git} e pastas de teste;
  \item \textbf{Compose}: declare \texttt{app}, \texttt{db} (postgres:18) e
  \texttt{redis}, com health checks e volume nomeado para o banco;
  \item \textbf{Migrações}: crie um serviço \texttt{migrate} (ou um entrypoint)
  que roda \texttt{prisma migrate deploy} antes de iniciar o app;
  \item \textbf{Seed}: execute \texttt{prisma db seed} na primeira subida
  (guardado por um marcador no banco ou script idempotente);
  \item \textbf{Validação}: rode \texttt{docker compose down -v} e depois
  \texttt{up -d --build} numa pasta limpa; o app deve subir completo;
  \item \textbf{Documentação}: escreva no \texttt{README} como subir o
  ambiente inteiro com um comando.
\end{passos}

\section{Exercícios de fixação}

\begin{exercicios}
  \item Qual a diferença entre imagem, contêiner e volume?
  \item Explique o que é multi-stage e qual o ganho prático.
  \item Escreva um \texttt{docker-compose.yml} para um app Node + PostgreSQL.
  \item O que faz um health check e por que \texttt{depends\_on} sozinho não basta?
  \item Liste 3 comandos Docker para inspecionar o que está rodando.
\end{exercicios}

\section{Desafios}

\begin{desafios}
  \item \textbf{Redução de imagem}: analise o tamanho das camadas com \texttt{docker history} e reduza a imagem final em pelo menos 30\%.
  \item \textbf{Dev vs prod}: crie perfis (\texttt{--profile}) que sobem serviços extras (adminer, mailhog) só em desenvolvimento.
  \item \textbf{Scan de segurança}: rode \texttt{docker scout} (ou trivy) na imagem e corrija as vulnerabilidades críticas.
\end{desafios}

\section{Exercícios de nível sênior}

\begin{seniores}
  \item \textbf{Orquestração}: compare Docker Compose com Kubernetes para produção e descreva quando cada um é apropriado.
  \item \textbf{Segredos}: configure a injeção de segredos (Docker secrets ou variáveis do CI) e explique por que a imagem nunca deve contê-los.
  \item \textbf{Cache de build}: otimize o cache do Docker (BuildKit) para que um rebuild de código leve $<$ 30s.
  \item \textbf{Observabilidade}: configure o \texttt{logging} do Compose para enviar logs ao Loki/ELK (base do capítulo~\ref{cap:observabilidade}).
\end{seniores}

\section{Armadilhas comuns}

\begin{itemize}[leftmargin=*, itemsep=0.4em]
  \item Imagens gigantes com ferramentas de build desnecessárias;
  \item Rodar como root dentro do contêiner;
  \item \texttt{latest} sem pinar versões (build não reproduzível);
  \item Esquecer volumes — dados somem a cada \texttt{down};
  \item Secrets gravados na imagem (vazam para qualquer um com acesso ao registry);
  \item \texttt{localhost} dentro do contêiner apontando para o host em vez de outro serviço;
  \item Não usar \texttt{.dockerignore} — o build copia \texttt{node\_modules}
  e \texttt{.git} para o contexto, deixando o build lento e a imagem inchada.
\end{itemize}

\section{Resumo do capítulo}

\begin{itemize}[leftmargin=*, itemsep=0.3em]
  \item Contêineres = aplicação + ambiente, leves e portáveis;
  \item Dockerfile multi-stage com usuário não-root é o padrão;
  \item Compose orquestra app + banco + cache com health checks;
  \item Volumes persistem; redes conectam; imagens pinadas reproduzem;
  \item Projeto entregue: SkillHub 100\% dockerizado.
\end{itemize}

\section{Referências oficiais para aprofundar}

\begin{itemize}[leftmargin=*, itemsep=0.3em]
  \item Documentação oficial do Docker: \url{https://docs.docker.com}
  \item Guia de boas práticas de Dockerfile: \url{https://docs.docker.com/develop/develop-images/dockerfile_best-practices/}
  \item Docker Compose: \url{https://docs.docker.com/compose/}
  \item Trivy (scan de vulnerabilidades): \url{https://trivy.dev}
\end{itemize}
